\section{Physical AI：Omniverse + Cosmos 的系统组合}

\subsection{从语言常识到物理常识}

Jensen 把 Physical AI 的难点说得很直接：机器人不仅要“会说”，还要“懂世界如何运动”。语言模型解决的是语义推断，机器人系统需要额外掌握物理约束、空间关系与执行反馈。

\begin{importantbox}{一句话定义 Physical AI}
\textit{``Everything that moves will be robotic someday.''}

这句话在工程上对应三件事：可感知（Perception）、可决策（Planning）、可执行（Control）。任何一环不稳定，系统就无法规模化落地。
\end{importantbox}

\subsection{为什么要先建“可扩展的虚拟训练场”}

访谈中提到，真实世界训练昂贵且慢，机器人硬件还会产生磨损。Omniverse 的价值在于提供可控、可重复、可并行的 3D 物理环境；Cosmos 的价值在于给这些环境提供更丰富的世界先验与场景生成能力。

\begin{knowledgebox}{Omniverse 与 Cosmos 的分工}
\begin{itemize}
    \item Omniverse：仿真引擎与数字孪生工作台，负责“物理一致性”和“训练执行”。
    \item Cosmos：世界模型能力，负责“场景多样性”和“语义-物理桥接”。
\end{itemize}
二者组合的目标是把机器人训练从“昂贵试错”转为“可工业化迭代”。
\end{knowledgebox}

\begin{warningbox}{误区：有了 world model 就能直接替代仿真}
world model 与仿真不是替代关系。前者擅长生成和泛化，后者擅长约束和验证。机器人系统需要“生成能力 + 物理一致性 + 闭环评估”三者并存，缺一不可。
\end{warningbox}

\subsection{本章小结}
Physical AI 的关键不是单点模型性能，而是训练基础设施。Omniverse + Cosmos 代表的是“机器人时代的数据工厂 + 先验工厂”。

